\chapter{Introduction}
\section{Nuclear physics: understanding the fundamental constituents of matter}
All ordinary matter is made of atoms. Named after the Greek word {\it atomos}, meaning ``indivisible," atoms were once thought to be the most basic building blocks of nature. Today, we know that atoms are composite structures built from electrons, protons, and neutrons. Furthermore, protons and neutrons, collectively called ``nucleons," are themselves composed of quarks and gluons. The electron, quarks, and gluons are examples of elementary particles, the truly irreducible {\it atomos} of the universe. They are members of the Standard Model of particle physics, which is currently the best physical theory capable of describing three of the four fundamental forces: electromagnetism, the strong force, and the weak force. (The fourth force, gravity, is described by the theory of general relativity.) 

Although the existence of quarks and gluons has been experimentally verified, many aspects of their behavior remain elusive. How are the quarks and gluons distributed inside of protons and neutrons? How do these extremely light particles generate nearly all the mass in the visible universe? How do their interactions give rise to both ordinary and exotic forms of matter, from everyday nucleons and nuclei to the primordial quark-gluon plasma (QGP) that existed in the early universe? These are examples of the fundamental questions we study in high-energy nuclear physics.

At the frontier of these pursuits today is the Electron-Ion Collider (EIC). Currently being built at Brookhaven National Lab on Long Island, the EIC will allow us to probe the quarks and gluons within nucleons with greater precision than previously possible. To support these future experiments, in this thesis we explore the application of machine learning towards the development of efficient data analysis for $D^0$-mesons produced during high-energy collisions, through which we can image the behavior of gluons inside the proton.

%----------------------------------------------------------
\section{Outline of this thesis}

This thesis was completed as part of the Jets and Heavy Flavor Working Group of the ePIC Collaboration \cite{ePIC}, which works on the development and evaluation of techniques for heavy flavor analysis at the future EIC \cite{EIC}. Codes and example results for our analyses can by found in the GitHub repository \texttt{ML\_D0\_analysis\_for\_ePIC} \cite{github}. These are adapted from the frameworks by Shyam Kumar \cite{Shyam_analysis,Shyam_ml} and Rongrong Ma \cite{Rongrong_jobs} located in the ePIC \texttt{snippets} repository.

We begin in Ch. \ref{ch.bkg} with background on the strong force and nuclear phenomena, the study of proton structure through deep-inelastic scattering, and the role of EIC experiments and heavy flavor probes. We then begin our analysis of $D^0$-mesons in Ch. \ref{ch:d0_top} with the examination of the $D^0\rightarrow\pi+K$ decay topology, observing the dependence of topological features on $D^0$ kinematics and motivating the use of topological cuts in the separation of $D^0$ signals from background. In Ch. \ref{ch.ml}, we introduce the machine learning concepts behind boosted decision trees and the aspects we consider in our BDT applications. Then in Ch. \ref{ch.results}, we describe and present results of our analyses. These include a preliminary analysis using an existing BDT framework to examine the impacts of machine-background effects on $D^0$ analysis, discussions of steps taken to improve the BDT methodology, and a comparison of our BDT methods to non-machine learning based topological cuts. 